\documentclass{article}
\usepackage[T1]{fontenc}
\usepackage{lmodern}
\usepackage{graphicx}
\usepackage{amsmath,amssymb}
\usepackage[a4paper,margin=2cm]{geometry}
\usepackage[absolute,overlay]{textpos}
\setlength{\TPHorizModule}{1mm}
\setlength{\TPVertModule}{1mm}
\usepackage[indonesian]{babel}
\usepackage{hyperref}
\setlength{\emergencystretch}{2em}
% unit: Halaman 11

\begin{document}

% === Halaman 11 (llm) ===

\begin{itemize}
    \item Garis besar Algoritma \textit{Rijndael} yang beroperasi pada blok 128-bit dengan kunci 128-bit adalah sebagai berikut (di luar proses pembangkitan \textit{round key}):
    \begin{enumerate}
        \item \textit{AddRoundKey}: melakukan XOR antara \textit{state} awal (plainteks) dengan \textit{cipher key}. Tahap ini disebut juga \textit{initial round}.
        \item Putaran sebanyak $Nr - 1$ kali. Proses yang dilakukan pada setiap putaran adalah:
        \begin{enumerate}
            \item \textit{SubBytes}: substitusi \textit{byte} dengan menggunakan tabel substitusi (S-box).
            \item \textit{ShiftRows}: pergeseran baris-baris \textit{array state} secara \textit{wrapping}.
            \item \textit{MixColumns}: mengacak data di masing-masing kolom \textit{array state}.
            \item \textit{AddRoundKey}: melakukan XOR antara \textit{state} sekarang \textit{round key}.
        \end{enumerate}
        \item \textit{Final round}: proses untuk putaran terakhir:
        \begin{enumerate}
            \item \textit{SubBytes}
            \item \textit{ShiftRows}
            \item \textit{AddRoundKey}
        \end{enumerate}
    \end{enumerate}
\end{itemize}

\end{document}
